\documentclass[../main.tex]{subfiles}

\begin{document}

% 使用方法：
% 1. 复制本文件并重命名为 files/2026.tex、files/2026f.tex 或 files/2026x.tex
% 2. 修改下面的元信息
% 3. 保持“题面章节 + 参考答案章节”的双章节结构
% 4. 在 main.tex 对应 part 中加入 \subfile{files/你的文件名}
%
% 约定总结：
% - 所有试卷文件都使用 subfiles，且以 \chapter 开始
% - 题号在不同 section 之间连续编号，依靠 \setcounter{enumi}{...} 控制
% - 普通题答案统一写成 \textbf{Solution}.，证明题统一写成 \textbf{Proof}.
% - 题面中的留白主要使用 \vspace{...}
% - 若某题含图，通常将图片文件放在 files/ 下，并命名为 <文件名>-figure*.pdf

% ====================
% 元信息：只改这里
% ====================
\newcommand{\ExamYearRange}{2025-2026}
\newcommand{\ExamCourse}{微积分（B）}
\newcommand{\ExamTerm}{（上）}
\newcommand{\ExamType}{期中考试}
\newcommand{\ExamTitle}{\ExamYearRange 学年\ExamCourse\ExamTerm\ExamType}

% ====================
% 常用写法
% ====================
\newcommand{\qblank}[1]{\vspace{#1}}
\newcommand{\sol}{\textbf{Solution}. }
\newcommand{\prf}{\textbf{Proof}. }

% ====================
% 模板 A：上学期 / 下学期期中
% 适用文件：2025.tex、2024.tex、2025x.tex 等
% ====================
\chapter{\ExamTitle}

\section{基本计算题(每小题 6 分，共 60 分)}
\begin{enumerate}
  \item 题目内容.
    \qblank{10em}

  \item 题目内容.
    \qblank{10em}
\end{enumerate}

\section{综合题(每小题 6 分，共 30 分)}
\begin{enumerate}
  \setcounter{enumi}{10}
  \item 题目内容.
    \qblank{12em}

  \item 题目内容.
    \qblank{12em}
\end{enumerate}

% 如果该卷含证明题，则保留本节；如果没有证明题，可整节删除。
\section{证明题(每小题 5 分，共 10 分)}
\begin{enumerate}
  \setcounter{enumi}{15}
  \item 题目内容.
    \qblank{10em}

  \item 题目内容.
    \qblank{10em}
\end{enumerate}

\chapter{\ExamTitle 参考答案}

\section{基本计算题(每小题 6 分，共 60 分)}
\begin{enumerate}
  \item \sol 参考答案内容.

  \item \sol 参考答案内容.
\end{enumerate}

\section{综合题(每小题 6 分，共 30 分)}
\begin{enumerate}
  \setcounter{enumi}{10}
  \item \sol 参考答案内容.

  \item \sol 参考答案内容.
\end{enumerate}

\section{证明题(每小题 5 分，共 10 分)}
\begin{enumerate}
  \setcounter{enumi}{15}
  \item \prf 证明内容.

  \item \prf 证明内容.
\end{enumerate}

% ====================
% 模板 B：上学期期末
% 使用时把上面的“模板 A”整体替换为下面结构即可
% ====================
%
% \chapter{\ExamTitle}
%
% \section{单项选择题(每小题 3 分，共 18 分)}
% \begin{enumerate}
%   \item 题目内容（\qquad）.
%
%   \begin{tasks}[label=\textbf{\Alph*.}, label-width=1.5em, column-sep=1em](2)
%     \task 选项A
%     \task 选项B
%     \task 选项C
%     \task 选项D
%   \end{tasks}
%   \qblank{1em}
% \end{enumerate}
%
% \section{填空题(每小题 4 分，共 16 分)}
% \begin{enumerate}
%   \setcounter{enumi}{6}
%   \item 题目内容$\underline{\hspace{3cm}}$.
%     \qblank{1em}
% \end{enumerate}
%
% \section{计算题(每小题 7 分，共 42 分)}
% \begin{enumerate}
%   \setcounter{enumi}{10}
%   \item 题目内容.
%     \qblank{10em}
% \end{enumerate}
%
% \section{综合题(每小题 7 分，共 14 分)}
% \begin{enumerate}
%   \setcounter{enumi}{16}
%   \item 题目内容.
%     \qblank{12em}
% \end{enumerate}
%
% \section{证明题(每小题 5 分，共 10 分)}
% \begin{enumerate}
%   \setcounter{enumi}{18}
%   \item 题目内容.
%     \qblank{15em}
% \end{enumerate}
%
% \chapter{\ExamTitle 参考答案}
%
% \section{单项选择题(每小题 3 分，共 18 分)}
% \begin{enumerate}
%   \item \sol A.
% \end{enumerate}
%
% \section{填空题(每小题 4 分，共 16 分)}
% \begin{enumerate}
%   \setcounter{enumi}{6}
%   \item \sol 答案内容.
% \end{enumerate}
%
% \section{计算题(每小题 7 分，共 42 分)}
% \begin{enumerate}
%   \setcounter{enumi}{10}
%   \item \sol 参考答案内容.
% \end{enumerate}
%
% \section{综合题(每小题 7 分，共 14 分)}
% \begin{enumerate}
%   \setcounter{enumi}{16}
%   \item \sol 参考答案内容.
% \end{enumerate}
%
% \section{证明题(每小题 5 分，共 10 分)}
% \begin{enumerate}
%   \setcounter{enumi}{18}
%   \item \prf 证明内容.
% \end{enumerate}

% ====================
% 模板 C：下学期期中（无证明题）
% 这一类通常是“基本计算题 + 综合题”
% 综合题常见为每小题 8 分，共 40 分
% ====================
%
% \chapter{\ExamTitle}
%
% \section{基本计算题(每小题 6 分，共 60 分)}
% \begin{enumerate}
%   \item 题目内容.
%     \qblank{10em}
% \end{enumerate}
%
% \section{综合题(每小题 8 分，共 40 分)}
% \begin{enumerate}
%   \setcounter{enumi}{10}
%   \item 题目内容.
%     \qblank{10em}
% \end{enumerate}
%
% \chapter{\ExamTitle 参考答案}
%
% \section{基本计算题(每小题 6 分，共 60 分)}
% \begin{enumerate}
%   \item \sol 参考答案内容.
% \end{enumerate}
%
% \section{综合题(每小题 8 分，共 40 分)}
% \begin{enumerate}
%   \setcounter{enumi}{10}
%   \item \sol 参考答案内容.
% \end{enumerate}

\end{document}
